\documentclass[11pt]{article}
\usepackage{icmlko}
\begin{document}

\title{고를 것이 없어서가 아니라 고를 눈이 없어서 --- 도메인마다 모형을 고르게 했더니 동전 던지기와 구별이 안 됐다}
\author{월드모델 실험실}
\date{노트 562 · 2026년 8월 5일}
\maketitle

\begin{abstract}
노트 $437$ 이 남긴 것 --- \textbf{이웃 효과(②)를 잡는 손잡이가 없다}.
\texttt{DOMDROP} 은 자기 축(①)만 잡는데, 크기는 ②가 더 크다: 도서에서 ②만
\textbf{$+0.128$}(전역 $+0.0459$ $-$ 자기 몫 $-0.0820$)이고 팝업은
$+0.0374$ 가 \emph{전부} ②다.

이 루프의 틀이 \textbf{포트폴리오}이므로 답은 자연스럽다 --- \textbf{도메인마다
모형을 고른다}. 배정은 학습 안쪽 시간 분할에서만 골랐다(유보를 한 번도 안
봤다).

\begin{center}
\begin{tabular}{lrl}
\hline
팔 & 판 차 & \\
\hline
B 항상 Q & $-0.0034$ & \\
\textbf{C 배정} & $\mathbf{-0.0017}$ & 안쪽 검증에서 고름 \\
\textbf{D 무작위} & $\mathbf{-0.0014}$ & 같은 Q 비율, 도메인은 무작위 \\
\hline
\end{tabular}
\end{center}

\textbf{$C - D = -0.0003$.} 고르기가 동전 던지기보다 \emph{못하다}.
\end{abstract}

\begin{figure}[t]
\centering
\includegraphics[width=\textwidth]{figs/f1.pdf}
\caption{\textbf{가} --- 무작위와 같다. \textbf{나} --- 검증이 너무 작다.
\textbf{다} --- 잡음이 효과보다 크다.}
\end{figure}

\section{대조 팔이 없었으면 채택할 뻔했다}

C 는 B 보다 낫다($-0.0017$ 대 $-0.0034$). 그것만 보면 ``고르기가
값을 한다''로 읽힌다. \textbf{그런데 D 가 같은 비율로 똑같이 잘한다.}
C 가 B 보다 나은 것은 \emph{Q 를 덜 섞어서}일 뿐이고, 어느 도메인에 섞느냐는
아무 상관이 없다.

\begin{quote}
사전 등록에 \textbf{``무작위 배정 팔을 같은 씨앗으로 돌린다''}고 미리 적어
둔 것이 값을 했다. 노트 $133$(첫 양수를 그대로 채택하지 않는다)을
\emph{팔로} 만들어 넣은 것이다.
\end{quote}

\section{왜 못 고르나 --- 눈이 없다}

\begin{center}
\begin{tabular}{lrrr}
\hline
도메인 & 안쪽 학습 & 안쪽 검증 & Q 를 고른 횟수 \\
\hline
웹툰 & $2{,}089$ & $849$ & $1/12$ \\
게임 & $159$ & $100$ & $\mathbf{12/12}$ \\
아이돌 & $77$ & $45$ & $6/12$ \\
\textbf{팝업} & $\mathbf{0}$ & $24$ & $9/12$ \\
\textbf{도서} & $69$ & $\mathbf{11}$ & $0/12$ \\
\hline
\end{tabular}
\end{center}

\textbf{팝업은 안쪽 학습이 $0$ 행}이다 --- 안쪽 모형에 그 도메인이 아예 없고,
집 밖으로 떨어진 예보의 $\rho$ 로 배정을 골랐다. 시장팝업은 $4$ 행으로
문턱 미만이다.

\begin{quote}
$\rho$ 의 표준오차는 대략 $1/\sqrt{n-3}$ 이다. 안쪽 검증이 $200$ 행이면
$0.071$ 이고, \textbf{P 와 Q 의 도메인별 차는 $|0.0004|$ $\sim$ $|0.055|$
다.} 차가 잡음에 통째로 묻힌다. $11$ 개 중 \textbf{$3$ 개만 만장일치}
(게임 $12/12$ · 모바일 $0/12$ · 도서 $0/12$)이고 나머지 여덟은 씨앗마다
뒤집힌다.
\end{quote}

\section{크기 법칙}

효과 $e$ 를 고르려면 표준오차가 그 절반보다 작아야 한다 ---
$1/\sqrt{n} < e/2$, 즉 $\mathbf{n > (2/e)^2}$.

\begin{center}
\begin{tabular}{lr}
\hline
효과 & 필요한 검증 행 \\
\hline
$0.05$ & $1{,}600$ \\
$0.02$ & $10{,}000$ \\
$0.005$ & $160{,}000$ \\
\hline
\end{tabular}
\end{center}

\textbf{도메인 $11$ 개 $\times$ $1{,}600$ 행 $=$ $17{,}600$ 행이 필요한데
안쪽 검증 전체가 $3{,}669$ 행이다 --- 다섯 배 부족하다.}

\begin{quote}
노트 $506$ 이 ``\textbf{자료가 유일한 배율}''이라 적은 것이 \emph{고르기
층에서} 다시 나온다. 모형을 더 잘 만드는 것이 병목이 아닌 것처럼,
\textbf{모형을 더 잘 고르는 것도 병목이 아니다} --- 고를 눈을 살 자료가 없다.
\end{quote}

\section*{노트 551 의 조항이 또 구했다}

안쪽 검증이 도서에 \textbf{$0/12$ 로 P 를} 골랐다. 노트 $559$ 의 전역 표에서
도서가 $+0.0459$ 로 제일 컸으니 ``제일 큰 이득을 놓쳤다''로 읽힐 자리다.

\begin{quote}
\textbf{그런데 도서만 바꾼 진짜 값은 $-0.0820$ 이다}(노트 $560$). 선택이
\emph{옳았다}. \textbf{전역 표로 선택을 채점하는 것도 같은 함정}이었고, 조항이
세 번째로 값을 했다.
\end{quote}

\section*{조항}

\begin{quote}
\textbf{고르는 단위마다 검증 표본이 효과보다 커야 고를 수 있다.} 최소 크기는
대략 $(2/e)^2$ 다. 그 아래면 \textbf{무작위 배정과 구별이 안 되고}, 그것을
확인하려면 \textbf{무작위 배정 팔을 같은 씨앗으로 돌려야 한다.}
\end{quote}

\section*{남은 것}

\begin{itemize}
\item \textbf{게임만 $12/12$} --- 안쪽 검증 $100$ 행인데도 만장일치다. 단일
전환으로 재 볼 수 있는 유일한 후보(무게 $5.3\%$).
\item \textbf{이웃 효과 ②를 잡는 손잡이는 여전히 없다.} \texttt{DOMDROP} 은
①만, 모형 고르기는 눈이 없다. 남은 길은 도메인별 가중이나 도메인$\times$축
상호작용인데 후자는 노트 $540$ 이 ``도메인 지시자''로 닫았다.
\item \textbf{안쪽 분할을 무작위로 하면} 검증이 커지지만 시간 유출이 생긴다 ---
이 판의 모든 채점이 시간 분할 위에 서 있으므로 규약을 바꾸는 일이다.
\end{itemize}

\end{document}
